\subsection{Process Mining and Object-Centric Event Logs}
\label{sec:bg-pm}

Process mining operates on event logs to discover process models, check
conformance, and enhance existing
models~\cite{van_der_aalst_process_mining_2011,van_der_aalst_2016_process_mining}.
The classical event log $L$ is a multiset of traces, each trace a sequence of
events over an activity alphabet; an event carries at minimum an activity label,
a timestamp, and a case identifier. Conformance checking relates a log to a
model either by token replay~\cite{rozinat_van_der_aalst_2008} or by computing
optimal alignments~\cite{adriansyah_et_al_alignments_2011,adriansyah_2014_phd},
yielding fitness and precision measures over the four quality
dimensions~\cite{buijs_et_al_quality_dimensions_2012}. The control-flow
semantics available to a declared model are systematised in the workflow
patterns literature~\cite{russell_workflow_patterns_2015}, and reference
implementations of the standard algorithm catalogue are provided by
PM4Py~\cite{berti_pm4py_2023}.

The single-case assumption of classical logs breaks down in enterprise settings
where one event touches many entities --- a payment references an invoice, a
purchase order, and a bank transaction simultaneously. \emph{Object-Centric
Event Logs} (OCEL)~\cite{ghahfarokhi_et_al_ocel_2021}, standardised as
OCEL~2.0~\cite{berti_et_al_ocel2_spec_2024,ocel_2_standard_2023}, resolve this:
an event is associated with a \emph{set} of typed objects, and the log is
formally a structure
$(\mathcal{E}, \mathcal{O}, \mathcal{A}, \mathcal{T},
\textit{evAttr}, \textit{objAttr}, \textit{E2O}, \textit{O2O})$
comprising events, objects, activities, object types, attribute maps, and the
event-to-object and object-to-object relations. Object-centric Petri
nets~\cite{van_der_aalst_berti_2020} and object-centric directly-follows
multigraphs~\cite{berti_van_der_aalst_oc_dfg_2023} give the discovery-side
semantics; OCPQ~\cite{kuesters_van_der_aalst_ocpq_2025} provides constraint
querying over these structures, and declarative object-centric constraints are
treated in~\cite{kuesters_van_der_aalst_oc_declare_2025}. Van der Aalst argues
that object-centric process intelligence is the prerequisite for trustworthy
enterprise AI~\cite{van_der_aalst_no_ai_without_pi_2024} --- a position this
paper operationalises by making the object-event relation a signed, first-class
component of every admitted change.

In our running P2P example, the \textsf{Pay} event is related to three objects
--- \textsf{Invoice\,\#4711}, \textsf{PO\,\#982}, \textsf{BankTx\,\#5530} ---
and the three-way-match rule is expressible only because the log model carries
all three relations on a single event. A classical single-case log would force
a choice of case notion and lose exactly the cross-object structure that the
fraud checks need.

\paragraph{Boundary predicates.}
The \emph{boundary predicate} $B : S \times E \to \{0,1\}$, where $S$ is the
set of process states and $E$ the set of event types, encodes whether an event
type is admissible from a given state. If $B(s,e) = 0$, the state transition
function must satisfy $\mu(s,e) = s$: the event produces no state change.
Boundary predicates appear implicitly throughout the conformance literature as
guard conditions on Petri net transitions and as the enabledness relation of
workflow nets~\cite{van_der_aalst_wfnet_soundness_1997}; recent work on
compliance-aware predictive monitoring uses them
prospectively~\cite{de_santis_compliance_predictive_2024}. We make the
predicate an explicit, versioned, first-class component of the admissibility
model: it is evaluated \emph{before} state mutation, its version is bound into
the change signature, and its refusals are enumerable
(Section~\ref{sec:impl-refusals}).

\subsection{Receipt Chains and the Axiom of Receipt Substance}
\label{sec:bg-receipts}

A \emph{receipt} $R_t$ is a cryptographic record witnessing that the state
transition $s_t \xrightarrow{e_t} s_{t+1}$ occurred with specific inputs and
outputs. We write $R_t \vdash s_{t+1} = \mu(s_t, e_t)$. Receipts compose into
a chain
\[
  R_0 \to R_1 \to \cdots \to R_n,
\]
where each $R_t$ includes the hash $h(R_{t-1})$ of its predecessor, the input
hash $h(\textit{in}_t)$, the output hash $h(\textit{out}_t)$, and the run
identifier. The chain hash function is BLAKE3, which offers 256-bit collision
resistance, structural immunity to length-extension attacks, and tree-hash
parallelism suitable for streaming over large event logs. Because hash-based
constructions rely only on preimage and collision resistance, the chain itself
is conservatively post-quantum even before signatures are considered; related
type-safe treatments of process evidence appear
in~\cite{type_safe_process_evidence_2024}, and provenance-chain constructions
over semantic-web data in~\cite{provenance_semantic_web_2025}.

A receipt is \emph{not} a receipt if it carries only hash values without the
backing event data. We adopt this as an axiom:

\begin{quote}
\itshape A receipt containing only expected hash, observed hash, alignment
state, and receipt hash is an unverifiable claim, not evidence.
\end{quote}

The axiom has teeth: it rules out an entire genre of audit theatre in which a
producer asserts ``hashes matched'' without the verifier being able to recompute
anything. A valid receipt in our framework must therefore embed, or
content-addressably reference, the full OCEL~2.0 logs it witnesses, so that any
holder of the receipt can re-canonicalise, re-hash, and re-align. The
consequence is formalised as Definition~\ref{def:valid-receipt}.

\paragraph{Proof classes.}
Receipts are ordered by evidential strength into a proof class hierarchy:
\[
  \textsc{FixtureProof} \prec \textsc{PredicateProof} \prec
  \textsc{ChicagoProof} \prec \textsc{BoundaryProof}.
\]
A \textsc{FixtureProof} witnesses only that a computation reproduced a stored
fixture; a \textsc{PredicateProof} adds machine-checked structural predicates;
a \textsc{ChicagoProof} requires end-to-end execution through a real runtime
boundary (no mocks at the verified interface); a \textsc{BoundaryProof} adds
runtime-observer envelopes and challenge nonces binding the evidence to a
specific physical execution. A receipt may not claim a class unless all
conditions of that class hold --- the refusal state
\texttt{ProofClassOverclaimed} (Section~\ref{sec:impl-refusals}) enforces the
hierarchy. The hierarchy matters for governance because it gives auditors a
vocabulary for \emph{how much} a given receipt proves, in the spirit of
capability maturity levels~\cite{sok_ccmm_2024}.

\subsection{Post-Quantum Cryptographic Signing}
\label{sec:bg-pqc}

Classical ECDSA and RSA signatures fall to Shor's algorithm on a
cryptographically relevant quantum computer~\cite{bernstein_lange_pqc_2017}.
This matters more for enterprise change receipts than for session
cryptography: a TLS session key protects data for minutes, but a receipt
signed today may need to carry evidential weight in a regulatory proceeding
fifteen years from now --- squarely inside the harvest-now--decrypt-later
window that has driven government migration
mandates~\cite{presman_govt_pqc_2024}. NIST's 2024 standards provide the
replacement primitives: FIPS~204 specifies \emph{ML-DSA}
(Module-Lattice-Based Digital Signature Algorithm, derived from
CRYSTALS-Dilithium)~\cite{nist_fips_204_2024}, and FIPS~205 specifies
\emph{SLH-DSA} (Stateless Hash-Based Digital Signature Algorithm, derived from
SPHINCS+)~\cite{nist_fips_205_2024}, whose security rests only on hash-function
assumptions and therefore serves as the conservative fallback should lattice
assumptions erode.

We model signing abstractly. The authority over a change computes
\[
  \sigma \;=\; \operatorname{Sign}\bigl(sk,\;
  h(a \,\|\, e \,\|\, O \,\|\, s \,\|\, R_{\mathrm{prior}} \,\|\, p \,\|\, v
  \,\|\, t \,\|\, n)\bigr),
\]
binding the actor $a$, event type $e$, object relations $O$, process state $s$,
prior receipt $R_{\mathrm{prior}}$, policy version $p$, validator version $v$,
timestamp $t$, and nonce $n$ under one signature. Verification is
$\verifysig(\sigma, pk, m) \in \{0,1\}$. The framework is algorithm-agnostic
--- any EUF-CMA-secure scheme suffices for the proofs of
Section~\ref{sec:inaccessibility} --- but the implementation default is
ML-DSA-65, and the deployment guidance assumes the crypto-agility posture of
\cite{crypto_agility_maturity_2022}: algorithm identifiers are versioned fields
of the receipt, so that a future migration re-signs the chain head rather than
rewriting history.
